\telubab{2}{BAB II}{LITERATURE REVIEW}{sec:2}

\begin{quote}
\textbf{DRAFT --- not submission-ready.} - \textbf{Venue status (checked
2026-10-11):} every venue in the reference list has been checked against
Semantic Scholar's publication record (batch lookup by arXiv ID),
Crossref, the official proceedings, or the paper itself. ``preprint''
means no peer-reviewed venue was found, so the paper is cited as arXiv.
Two items remain to confirm on a primary page: CaRR at ACL 2026
(Semantic Scholar only) and Search-R1 at COLM 2025 (the paper's own
statement; Semantic Scholar lists arXiv). - \textbf{Paper descriptions:}
these are drawn from abstracts and the project's research notes
(\texttt{docs/\allowbreak{}research/\allowbreak{}}). Before submission,
confirm each characterization against the full paper.
\end{quote}

This chapter reviews the background needed for the thesis:

\begin{itemize}
\tightlist
\item
  \textbf{Sections 2.1--2.4} cover the building blocks: language models,
  retrieval, multi-hop QA with evidence, and LLM agents.
\item
  \textbf{Section~\ref{sec:2.5}} covers reinforcement learning for LLMs,
  from PPO to GRPO and its recent fixes.
\item
  \textbf{Sections 2.6--2.7} review RL-trained search agents and
  evidence-based rewards, the area this thesis contributes to.
\item
  \textbf{Section~\ref{sec:2.8}} covers reward hacking and how RL
  results should be evaluated.
\item
  \textbf{Section~\ref{sec:2.9}} draws out the gaps the thesis
  addresses.
\end{itemize}

\subsection{Large language models and instruction tuning}\label{sec:2.1}

Large language models (LLMs) are first pretrained on next-token
prediction over large text corpora, then adapted to follow instructions
through supervised fine-tuning and preference-based training. Ouyang et
al.~(2022) showed that fine-tuning with human feedback makes models
follow user intent better, even when they are much smaller than the
original model.

This thesis uses \textbf{Qwen2.5-7B-Instruct} (Qwen Team, 2024,
technical report), an instruction-tuned open-weight model. A 7B model is
small enough to fine-tune on a single 48 GB GPU, as long as the training
is parameter-efficient (Section 2.5.4). The instruction tuning also
matters for the experiments: the untrained model can already attempt the
action format, which gives reinforcement learning a starting point.

\subsection{Retrieval-augmented generation and BM25}\label{sec:2.2}

An LLM's knowledge is fixed at training time and can be wrong.
\textbf{Retrieval-augmented generation (RAG)} addresses this by
retrieving relevant documents and conditioning the answer on them. In
basic RAG, retrieval happens once, before generation. Questions that
need several pieces of information may need several retrievals, each one
depending on what the previous one found. That motivates the agentic
approach in Section~\ref{sec:2.4}.

This thesis uses \textbf{BM25} (Robertson \& Zaragoza, 2009) as the
retriever. BM25 is a sparse lexical ranking function: it scores a
document by how often the query terms appear in it, weighted by how rare
each term is, and normalized by document length. It needs no training
and is deterministic, so a given query always returns the same
documents. This determinism is what makes the evidence reward in
Chapter~\ref{sec:3} exactly verifiable: it is always known which query
surfaced which document.

\subsection{Multi-hop QA, HotpotQA, and citation
evaluation}\label{sec:2.3}

\textbf{HotpotQA} (Yang et al., 2018) is a multi-hop question answering
dataset built from Wikipedia: answering a question requires combining
information from two paragraphs. It has two features that this thesis
relies on:

\begin{itemize}
\tightlist
\item
  \textbf{Distractor setting.} Each question comes with 10 paragraphs:
  the 2 gold paragraphs plus 8 distractors retrieved to look relevant.
  The system has to find the gold paragraphs among the distractors.
\item
  \textbf{Supporting facts.} Annotators marked the exact sentences, as
  (paragraph title, sentence index) pairs, that are needed to answer
  each question. The dataset evaluates answers with exact match (EM) and
  token F1, and evaluates supporting-fact prediction separately.
\end{itemize}

These sentence-level labels make it possible to check whether a model's
evidence is correct with a simple rule, without a human or a model
acting as judge.

\textbf{Citation evaluation.} Gao et al.~(2023 ✓, EMNLP) introduced
\textbf{ALCE}, a benchmark for LLMs that generate text with citations.
ALCE treats correctness and citation quality as separate dimensions, and
separates \emph{citation recall} (is every statement supported by the
cited passages?) from \emph{citation precision} (is every citation
relevant?). The authors report that even the best models of the time
often lacked complete citation support. This thesis uses the same split
between \emph{evidence coverage} (a form of recall) and \emph{citation
precision} (Chapter~\ref{sec:3}). The main difference is the setting:
here citations are checked against HotpotQA's gold sentences rather than
judged by a natural-language-inference model.

\subsection{LLM agents and tool use}\label{sec:2.4}

An \textbf{LLM agent} interleaves reasoning with actions in an
environment. \textbf{ReAct} (Yao et al., 2023) prompts a model to
alternate between reasoning and tool calls (such as a search) and to
condition each step on what the earlier steps observed. The agent in
this thesis follows a simplified ReAct-style loop with a strict action
grammar: on each turn it either searches or gives a final answer with
citations (Chapter~\ref{sec:3}).

Prompting alone does not guarantee that a model uses tools well. Jin et
al.~(2025 ✓, COLM) argue that prompted LLMs may not know how to interact
with a search engine optimally. This motivates training the agent's tool
use with reinforcement learning.

\subsection{Reinforcement learning for LLMs}\label{sec:2.5}

\subsubsection{Policy gradients and PPO}\label{sec:2.5.1}

In RL fine-tuning, the LLM is the \textbf{policy}: given a context, it
produces actions (tokens), and an episode is scored by a reward.
Policy-gradient methods raise the probability of actions that led to
above-average reward. \textbf{Proximal Policy Optimization} (PPO;
Schulman et al., 2017, preprint) stabilizes this by clipping the ratio
between the new and old policy probabilities, so that a single update
cannot move the policy too far. Standard PPO for LLMs also trains a
value network (critic) to estimate the expected reward, which roughly
doubles the memory needed.

\subsubsection{GRPO}\label{sec:2.5.2}

\textbf{Group Relative Policy Optimization} (GRPO) was introduced in
DeepSeekMath (Shao et al., 2024, preprint). It removes the critic.
Instead, it samples a group of outputs for the same prompt and uses each
output's reward relative to the group, normalized by the group's mean
and standard deviation, as its advantage.

\textbf{DeepSeek-R1} (DeepSeek-AI, 2025 ✓, Nature) showed that GRPO with
rule-based, verifiable rewards can train strong reasoning. Its rewards
checked answer accuracy and output format, with no learned reward model.
VEGR follows the same principle: every reward term can be checked by
rule.

Group normalization has a side effect that matters in this thesis. When
all outputs in a group receive the same reward, every advantage is zero,
and the group provides no learning signal.

\subsubsection{Fixes to GRPO: Dr.~GRPO and DAPO}\label{sec:2.5.3}

Two lines of recent work modify GRPO, and this thesis tests both.

\begin{itemize}
\tightlist
\item
  \textbf{Dr.~GRPO} (Liu et al., 2025, preprint) identifies an
  optimization bias in GRPO and removes its length and
  standard-deviation normalization terms. This thesis uses only the
  second change, computing advantages without dividing by the standard
  deviation (run R4).
\item
  \textbf{DAPO} (Yu et al., 2025 ✓, NeurIPS) proposes several changes
  for large-scale RL. One of them is \textbf{dynamic sampling}: groups
  whose outputs all get the same reward are dropped, and more groups are
  sampled until the batch has a learning signal. The thesis adapts this
  idea in run R0-dyn.
\end{itemize}

\subsubsection{Parameter-efficient fine-tuning: LoRA}\label{sec:2.5.4}

Fine-tuning all 7 billion parameters with RL does not fit on one GPU.
\textbf{LoRA} (Hu et al., 2022) freezes the pretrained weights and
trains small low-rank matrices added to selected layers. This cuts
memory to a fraction of full fine-tuning, which is what lets this
project train on a single 48 GB GPU.

\subsection{Reinforcement learning for search agents}\label{sec:2.6}

\textbf{Search-R1} (Jin et al., 2025 ✓, COLM) is the main reference
point for this thesis. It trains an LLM to issue search queries during
step-by-step reasoning, using RL with a \textbf{simple outcome-based
reward}: whether the final answer is correct. The training covers
multi-turn trajectories, and retrieved tokens are masked from the loss.
Search-R1 reports large gains over RAG baselines on seven QA datasets
with Qwen2.5 models. The outcome-only baseline in this thesis (R0)
follows the same reward design.

\textbf{AutoRefine} (Shi et al., 2025, NeurIPS) adds a step that refines
retrieved knowledge between searches. It trains with GRPO using both
answer-correctness and \textbf{retrieval-specific rewards}, and reports
its largest gains on multi-hop questions. This supports the idea, used
in VEGR, that rewarding intermediate search quality helps beyond the
final answer alone.

These systems share two properties that this thesis departs from:

\begin{enumerate}
\def\labelenumi{\arabic{enumi}.}
\tightlist
\item
  They are usually trained and reported at multi-GPU scale.
\item
  Their main reward is answer correctness. A correct answer does not
  show that the model used, or could point to, the evidence behind it.
\end{enumerate}

\subsection{Grounding and evidence-based rewards}\label{sec:2.7}

Several works move beyond single holistic rewards.

\textbf{Fine-Grained RLHF} (Wu et al., 2023) trains with several reward
models, each targeting a different kind of error (for example factual
incorrectness, irrelevance, or incompleteness), instead of one holistic
preference score. The authors show that changing how the rewards are
combined changes the model's behaviour. This supports VEGR's per-aspect
design and the re-weighting experiment R3.

Two 2026 papers address evidence directly in search-agent RL:

\begin{itemize}
\tightlist
\item
  \textbf{CaRR / C-GRPO} (Zhang et al., 2026, ACL) breaks multi-hop
  questions into single-hop rubrics. A judge LLM checks whether each
  rubric is supported by the cited content, and the rubric reward counts
  only when the answer is correct. The authors report that under
  outcome-only GRPO, agents learn \textbf{shortcut answers} and make
  fewer tool calls. This matches the shortcut behaviour reported in
  Chapter~\ref{sec:4}.
\item
  \textbf{STAMP} (Xu et al., 2026, preprint) credits each search step
  that first surfaced a cited supporting document, and uses an LLM
  verifier to check citations. It notes that citation-aware and
  evidence-coverage rewards are now common in this line of work.
\end{itemize}

\textbf{Position of this thesis.} Both CaRR and STAMP rely on an
\textbf{LLM judge or verifier}, which adds cost and possible bias. They
also work with live web search at larger model scales. VEGR instead uses
HotpotQA's gold supporting facts and a deterministic retriever, so every
evidence term is computed exactly by rule. The individual terms of VEGR
are not new (Section~\ref{sec:3.4}). What this thesis adds is a
judge-free combination for a small, single-GPU setting, together with a
controlled multi-seed test of what that combination changes.

\subsection{Reward hacking and evaluating RL}\label{sec:2.8}

\textbf{Reward hacking.} Skalse et al.~(2022 ✓, NeurIPS; the proceedings
title is ``Defining and Characterizing Reward Gaming'', while the arXiv
title is ``\ldots Reward Hacking'') formally define reward hacking:
optimizing an imperfect proxy reward lowers performance on the true
objective. They argue that proxies that cannot be hacked are rare. Gao,
Schulman and Hilton (2023) measure the same effect with learned reward
models: optimizing a proxy too hard eventually hurts the true (``gold'')
score. In this thesis, the proxy is answer correctness and the true
objective is a correct \emph{and} grounded answer. Chapter~\ref{sec:4}
reports cases where the outcome-only policy satisfied the proxy without
citing evidence.

\textbf{Evaluating RL with few runs.} RL results vary considerably
between random seeds.

\begin{itemize}
\tightlist
\item
  Henderson et al.~(2018) showed that differences between seeds alone
  can look like significant differences between methods.
\item
  Agarwal et al.~(2021) argue that when there are only a few runs,
  results should be reported as interval estimates (for example
  bootstrap confidence intervals), never as single runs.
\end{itemize}

This thesis therefore trains every main method with 3 seeds and compares
models with paired bootstrap confidence intervals
(Section~\ref{sec:3.7}).

\subsection{Research gaps and positioning}\label{sec:2.9}

{\def\LTcaptype{none} % do not increment counter
\begin{longtable}[]{@{}
  >{\raggedright\arraybackslash}p{(\linewidth - 4\tabcolsep) * \real{0.3333}}
  >{\raggedright\arraybackslash}p{(\linewidth - 4\tabcolsep) * \real{0.3333}}
  >{\raggedright\arraybackslash}p{(\linewidth - 4\tabcolsep) * \real{0.3333}}@{}}
\toprule\noalign{}
\begin{minipage}[b]{\linewidth}\raggedright
Gap
\end{minipage} & \begin{minipage}[b]{\linewidth}\raggedright
Evidence from literature
\end{minipage} & \begin{minipage}[b]{\linewidth}\raggedright
How this thesis addresses it
\end{minipage} \\
\midrule\noalign{}
\endhead
\bottomrule\noalign{}
\endlastfoot
\textbf{G1} Correct ≠ grounded & Outcome-only rewards (Search-R1) do not
check evidence; CaRR reports shortcut answers; reward hacking theory
(Skalse) & Measure grounding directly with sentence-level gold evidence;
compare R0 with VEGR \\
\textbf{G2} Evidence rewards need an LLM judge & CaRR and STAMP use LLM
judges or verifiers; ALCE uses an NLI model & VEGR checks evidence by
rule against HotpotQA supporting facts \\
\textbf{G3} Search-agent RL is reported at multi-GPU scale & Search-R1,
AutoRefine, CaRR, STAMP & Single 48 GB GPU, LoRA, about \$2.25 per
run \\
\textbf{G4} Single-run results mistake seed noise for method effects &
Henderson 2018; Agarwal 2021 & 3 seeds per method; paired bootstrap
CIs \\
\end{longtable}
}

In short, this thesis does not propose a new RL algorithm or claim
state-of-the-art accuracy. It asks a narrower question in a fully
controlled setting: \textbf{when the evidence can be checked exactly,
does rewarding it change how a search agent behaves, compared with
rewarding the answer alone?}

\begin{center}\rule{0.5\linewidth}{0.5pt}\end{center}

\subsection{References}\label{references}

Draft list. The ``Checked via'' column records how each venue was
verified on 2026-10-11. Final formatting follows the university's
citation style {[}TODO: which style{]}.

{\def\LTcaptype{none} % do not increment counter
\begin{longtable}[]{@{}
  >{\raggedright\arraybackslash}p{(\linewidth - 4\tabcolsep) * \real{0.3333}}
  >{\raggedright\arraybackslash}p{(\linewidth - 4\tabcolsep) * \real{0.3333}}
  >{\raggedright\arraybackslash}p{(\linewidth - 4\tabcolsep) * \real{0.3333}}@{}}
\toprule\noalign{}
\begin{minipage}[b]{\linewidth}\raggedright
Reference
\end{minipage} & \begin{minipage}[b]{\linewidth}\raggedright
Venue
\end{minipage} & \begin{minipage}[b]{\linewidth}\raggedright
Checked via
\end{minipage} \\
\midrule\noalign{}
\endhead
\bottomrule\noalign{}
\endlastfoot
Agarwal, R., Schwarzer, M., Castro, P. S., Courville, A., Bellemare, M.
G. (2021). Deep Reinforcement Learning at the Edge of the Statistical
Precipice. arXiv:2108.13264 & NeurIPS 2021 & Semantic Scholar \\
DeepSeek-AI (2025). DeepSeek-R1 & \emph{Nature} 2025 & earlier session
(writer brief) \\
Gao, L., Schulman, J., Hilton, J. (2023). Scaling Laws for Reward Model
Overoptimization. arXiv:2210.10760 & ICML 2023 & Semantic Scholar \\
Gao, T., Yen, H., Yu, J., Chen, D. (2023). Enabling Large Language
Models to Generate Text with Citations. pp.~6465--6488. arXiv:2305.14627
& EMNLP 2023 & ACL Anthology \\
Henderson, P., et al.~(2018). Deep Reinforcement Learning that Matters.
doi:10.1609/aaai.v32i1.11694 & AAAI 2018 & Semantic Scholar (AAAI
DOI) \\
Hu, E. J., et al.~(2022). LoRA: Low-Rank Adaptation of Large Language
Models. arXiv:2106.09685 & ICLR 2022 & Semantic Scholar \\
Jin, B., et al.~(2025). Search-R1: Training LLMs to Reason and Leverage
Search Engines with Reinforcement Learning. arXiv:2503.09516 & COLM 2025
& the paper's own statement {[}TODO: confirm on the COLM site{]} \\
Liu, Z., et al.~(2025). Understanding R1-Zero-Like Training: A Critical
Perspective. arXiv:2503.20783 & preprint & Semantic Scholar lists arXiv
only \\
Ouyang, L., et al.~(2022). Training language models to follow
instructions with human feedback. arXiv:2203.02155 & NeurIPS 2022 &
Semantic Scholar (proceedings DOI) \\
Yang, A., et al.~(Qwen Team) (2024). Qwen2.5 Technical Report.
arXiv:2412.15115 & technical report & Semantic Scholar \\
Robertson, S., Zaragoza, H. (2009). The Probabilistic Relevance
Framework: BM25 and Beyond. doi:10.1561/1500000019 & \emph{Foundations
and Trends in Information Retrieval} & Crossref \\
Schulman, J., et al.~(2017). Proximal Policy Optimization Algorithms.
arXiv:1707.06347 & preprint & --- \\
Shao, Z., et al.~(2024). DeepSeekMath: Pushing the Limits of
Mathematical Reasoning in Open Language Models. arXiv:2402.03300 &
preprint & Semantic Scholar lists arXiv only \\
Shi, Y., et al.~(2025). Search and Refine During Think: Facilitating
Knowledge Refinement for Improved Retrieval-Augmented Reasoning
(AutoRefine). arXiv:2505.11277 & NeurIPS 2025 & Semantic Scholar
(proceedings DOI 10.52202/085713-5211) \\
Skalse, J., Howe, N., Krasheninnikov, D., Krueger, D. (2022). Defining
and Characterizing Reward Gaming. arXiv:2209.13085 & NeurIPS 2022 &
NeurIPS proceedings \\
Xu, K., et al.~(2026). STAMP: Provenance-Guided Credit Assignment for
Deep Search Agents. arXiv:2607.11172 & preprint & Semantic Scholar \\
Wu, Z., et al.~(2023). Fine-Grained Human Feedback Gives Better Rewards
for Language Model Training. arXiv:2306.01693 & NeurIPS 2023 & Semantic
Scholar \\
Yang, Z., et al.~(2018). HotpotQA: A Dataset for Diverse, Explainable
Multi-hop Question Answering. doi:10.18653/v1/D18-1259 & EMNLP 2018 &
Semantic Scholar (ACL Anthology DOI) \\
Yao, S., et al.~(2023). ReAct: Synergizing Reasoning and Acting in
Language Models. arXiv:2210.03629 & ICLR 2023 & Semantic Scholar \\
Yu, Q., et al.~(2025). DAPO: An Open-Source LLM Reinforcement Learning
System at Scale. arXiv:2503.14476 & NeurIPS 2025 & NeurIPS
proceedings \\
Zhang, J., Lv, X., Feng, L., Hou, L., Li, J. (2026). Chaining the
Evidence: Robust Reinforcement Learning for Deep Search Agents with
Citation-Aware Rubric Rewards (CaRR). arXiv:2601.06021 & ACL 2026 &
Semantic Scholar only {[}TODO: confirm in the ACL Anthology{]} \\
\end{longtable}
}
